\section{A sequence of linear forms}\label{sec:forms}
\subsection{Zeros, poles, and the choice of parameters}
The integer zeros have two tasks: they remove initial terms from the
infinite sum and give identities for the rational correction terms.
The half-integer poles produce beta values. Varying their orders
across the pole interval improves the balance between denominator
growth and analytic decay.

Two differentiations raise each partial-fraction exponent by two.
Reflection will remove the odd beta indices, so a largest pole order
of $31$ leaves $32$ as the largest even index. The other parameters
below are chosen so that the decay exceeds the denominator cost,
as verified in Section~\ref{sec:decay}.

Fix
\begin{equation}\label{eq:parameters}
 S=31,\qquad r=\frac34,\qquad N=96,\qquad
 \delta_j=\frac{j}{96}\quad(0\le j\le30),
\end{equation}
and let $n$ range over positive multiples of $N$. Define the blocks
and their combined pole orders by
\[
 I_j(n)=[\delta_jn,(1-\delta_j)n]\cap\Z,\qquad
 B_a(n)=\#\{j:a\in I_j(n)\}\quad(0\le a\le n).
\]
The block $I_0(n)$ contains every address. At $n=96$, the successive
blocks are $[0,96],[1,95],\ldots,[30,66]$, so the orders rise from
one at the ends to $31$ in the middle. The base $N=96$ makes these
endpoints and $rn$ integral. We often suppress $n$ in $B_a(n)$.

Put
\[
 \Lambda_n=[-(1+r)n-1,rn]\cap\Z,
 \qquad
 R_n(t)=\frac{(2t+n+1)\prod_{\ell\in\Lambda_n}(t-\ell)^3}
 {\prod_{a=0}^{n}(t+a+\tfrac12)^{B_a}}.
\]
Two degree parameters will track the factorial normalization:
\begin{equation}\label{eq:degree-parameters}
 \gamma=\sum_{j=0}^{30}(1-2\delta_j)=\frac{341}{16},
 \qquad G=\gamma-3(1+2r)=\frac{221}{16}.
\end{equation}
There are $\gamma n+31$ denominator factors and
$3\#\Lambda_n=\tfrac{15}{2}n+6$ integer zero factors, counted with
multiplicity. Including the linear prefactor gives
$\deg R_n=-Gn-24$, so $R_n$ is proper and decays rapidly at infinity.
The prefactor cancels one order at the central pole; nominal pole
orders remain convenient, with the corresponding leading coefficient
equal to zero.

Reflection about $-(n+1)/2$ permutes the zeros and the pole blocks.
Since $\gamma n$ and $3\#\Lambda_n$ are even, the prefactor and
denominator both change sign. Hence
\begin{equation}\label{eq:reflection}
 R_n(-n-1-t)=R_n(t).
\end{equation}

\subsection{The beta expansion}
Write the partial fractions as
\[
 R_n(t)=\sum_{a=0}^{n}\sum_{i=1}^{B_a}
                   \frac{a_{i,a}}{(t+a+\tfrac12)^i},
\]
and set absent coefficients to zero. Let
\[
 (i)_2=i(i+1),\qquad
 \sigma_a(s)=\sum_{h=0}^{a-1}\frac{(-1)^h}{(2h+1)^s}.
\]
Differentiating twice and shifting each beta series gives
\begin{align}
 L_n&=\sum_{t=0}^{\infty}(-1)^tR_n''(t)
       =\rho_n+\sum_{i=1}^{31}\tau_{i,n}\beta(i+2),
                                                \label{eq:forms}\\
 \tau_{i,n}&=2^{i+2}(i)_2\sum_{a=0}^{n}(-1)^a a_{i,a},
                                                \label{eq:beta-coeff}\\
 \rho_n&=-\sum_{a=0}^{n}\sum_{i=1}^{B_a}
              (-1)^a2^{i+2}(i)_2a_{i,a}\sigma_a(i+2).
                                                \label{eq:constant}
\end{align}
All exponents in the shifted series are at least three, so the
rearrangements are absolutely convergent. Uniqueness of partial
fractions in \eqref{eq:reflection} gives
$a_{i,n-a}=(-1)^ia_{i,a}$. Since $n$ is even, every $\tau_{i,n}$
with odd $i$ vanishes. Thus
\begin{equation}\label{eq:target-span}
 L_n\in\Q+\Q\beta(4)+\Q\beta(6)+\cdots+\Q\beta(32).
\end{equation}

The triple zeros give two identities we shall use separately:
\begin{equation}\label{eq:zeros}
 R_n''(t)=0\quad(0\le t\le rn),\qquad
 R_n''(-\ell)=0\quad(1\le\ell\le n),
\end{equation}
where the arguments are integers. The first starts the analytic sum
beyond $rn$. The second controls the denominators of $\rho_n$.
